\bookchapter{第二十五章：二手手機處理與數字遺產}

\chapterintro{嘿，夥計們，經過前面二十多章的“修煉”，相信你已經從一個“數字小白”成長為一名相當從容的手機使用者了。但科技生活總有輪迴——新歡終會變舊愛。當你準備讓老夥計“退役”時，你知道該如何體面、安全地告別嗎？}

\booksection{一、舊手機的“安樂死”：安全出售或轉贈指南}

核心原則：資料清理不是“刪除”，而是“覆蓋”。許多人以為把照片拖進回收站、清空相簿就萬事大吉，結果手機轉手後，新主人用資料恢復軟體一掃，老照片、聊天記錄、身份證照片全回來了——這種“社死現場”每年都在上演。

\begin{enumerate}
  \item \textbf{自我"大掃除"}：完成最後一次完整備份，退出所有帳號（iCloud、華為/小米/三星帳號、Google帳號），關閉"查詢我的手機"功能——否則下一位使用者會被啟用鎖攔在門外。
  \item \textbf{執行終極"格式化"}：現代手機的"恢復出廠設定"通常包含加密擦除——它銷燬的是加密金鑰本身，金鑰一沒，整盤資料瞬間變為亂碼。
  \item \textbf{最後的"障眼法"（僅老裝置需要）}：現代手機（iOS 8+、Android 6+）已預設全盤加密（FBE），恢復出廠設定會銷燬加密金鑰，資料理論上不可恢復，\textbf{無需此步驟}。僅對極老舊的非加密裝置（2015年前），才需要在恢復出廠設定後用垃圾檔案填滿儲存，再執行一次恢復出廠設定。
\end{enumerate}

\textbf{各品牌資料擦除路徑速查表}：儘管原理相通，但每個品牌把“恢復出廠設定”藏的位置略有不同，關鍵差異如下。

\begin{center}
\begin{tabularx}{\linewidth}{|X|X|X|}
\hline
\textbf{品牌} & \textbf{擦除入口} & \textbf{關鍵差異與注意事項} \tabularnewline
\hline
華為 & 設定 > 系統和更新 > 重置 > 恢復出廠設定 & 需先退出華為帳號並關閉"查詢我的手機"，否則會觸發啟用鎖；純淨模式下可能需要密碼驗證。 \tabularnewline
小米 / Redmi & 設定 > 我的裝置 > 恢復出廠設定 & 底部有"格式化虛擬SD卡"和"清除全部內容與設定"選項，二者都要勾選；退出小米帳號前先關閉"查詢裝置"。 \tabularnewline
OPPO / 一加 & 設定 > 其他設定 > 還原手機 > 徹底清除全部資料 & ColorOS 提供"徹底清除"模式，會多次覆蓋儲存；需輸入鎖屏密碼確認。 \tabularnewline
vivo & 設定 > 系統管理 > 備份與重置 > 清除所有資料 & 勾選"格式化手機儲存"才會清理SD卡區；需退出vivo帳號，關閉"查詢手機"。 \tabularnewline
三星 & 設定 > 通用管理 > 重置 > 恢復出廠設定 & 需先移除Samsung帳號並關閉"查詢我的手機"；部分機型有"自動覆蓋"選項，建議開啟。 \tabularnewline
iPhone & 設定 > 通用 > 傳輸或還原iPhone > 抹掉所有內容和設定 & 必須先退出Apple ID，否則下一位使用者被iCloud啟用鎖攔住；iOS 15+支援"退出時保留iCloud資料"以備後查。 \tabularnewline
\hline
\end{tabularx}
\end{center}

\begin{tcbwarn}{最容易被忽略的三件事}
\begin{itemize}
  \item \textbf{未退出帳號}：手機賣給下家後被遠端鎖定，自己也下不來，最後只能持身份證去官方售後解鎖。
  \item \textbf{SD卡未取出}：恢復出廠設定不會格式化外接SD卡裡的照片，記得手動取出或單獨格式化。
  \item \textbf{隨身碟/OTG外設未斷開}：擦除過程中如果連著外設，可能誤清外設資料，也可能導致擦除中斷。
\end{itemize}

\end{tcbwarn}

\booksection{二、數字遺產：你的社交靈魂如何安放？}

手機裡的資料，早就不是“虛擬”的了。它承載著你的資金、社交關係、創作成果，乃至你的“數字身份”。\textbf{一個人去世後，這些資產該怎麼處理？}這不再是一個抽象問題，而是越來越多家庭真實面臨的難題——親屬不知道密碼、平臺不讓繼承、加密貨幣永遠丟失。

\begin{tcbinfo}{法律依據：個人資訊保護法第49條}
2021年施行的《中華人民共和國個人資訊保護法》第49條明確規定：\textbf{自然人死亡的，其近親屬為了自身的合法、正當利益，可以對死者相關個人資訊行使本章規定的查閱、複製、更正、刪除等權利；死者生前另有安排的除外。}

這意味著：

\begin{itemize}
  \item \textbf{近親屬有法律依據}向平臺申請處理死者的帳號與個人資訊。
  \item "合法、正當利益"通常包括辦理遺產繼承、關閉帳號、刪除敏感隱私等。
  \item 死者生前若有遺囑或協議另行安排（如指定遺產聯絡人），從其約定。
  \item 平臺不能以"使用者隱私"為由一概拒絕配合，但具體流程各家不同。
\end{itemize}

\end{tcbinfo}

\booksubsection{微信與支付寶：帳號不能繼承，但資金可以}

這兩個中國人最常用的App，規則各不相同：

\begin{itemize}
  \item \textbf{微信帳號}：依據《騰訊微信軟體許可及服務協議》，帳號所有權歸騰訊，使用者僅享有使用權，\textbf{帳號本身不能繼承、不能轉讓}。但帳號內的\textbf{零錢、零錢通資金}屬於個人財產，近親屬可憑死亡證明、親屬關係證明、繼承公證書向財付通申請提現或轉移。
  \item \textbf{支付寶}：帳號同樣不能直接繼承，但支付寶餘額、餘額寶、理財通等\textbf{資金屬於可繼承財產}。繼承人可透過支付寶"客服熱線"或線下提交死亡證明、繼承公證書辦理；螞蟻財富類理財產品需走基金公司流程。
  \item \textbf{聊天記錄}：法律上屬個人隱私，平臺不會主動提供給家屬。如涉及遺產線索，可透過訴訟或公證申請調取。
\end{itemize}

\begin{tcbwarn}{千萬別做的事}
\begin{itemize}
  \item \textbf{不要把死者帳號密碼直接交給家人"接著用"}——這違反平臺協議，資金被凍結甚至觸發風控時舉證更難。
  \item \textbf{不要嘗試用死者手機收驗證碼"提現"}——可能被認定為盜用，反而觸犯法律。
  \item \textbf{不要拖延}：部分帳號長期未登入會被自動登出，資金找回會更麻煩。
\end{itemize}

\end{tcbwarn}

\booksubsection{Apple“遺產聯絡人”：提前指定，事半而功倍}

iOS 15.2 起蘋果推出“遺產聯絡人”功能，允許使用者提前指定信任的聯絡人，去世後該聯絡人可憑\textbf{死亡證明 + 遺產聯絡人訪問金鑰}申請訪問死者的iCloud資料（照片、備忘錄、通訊錄、備忘錄等）。操作步驟：

\begin{enumerate}
  \item 在iPhone上進入\textbf{[設定] > [你的姓名] > [登入與安全] > [遺產聯絡人]}。
  \item 點選"新增遺產聯絡人"，使用面容ID確認身份。
  \item 選擇聯絡人（從通訊錄挑選），系統會生成一份\textbf{訪問金鑰}。
  \item 將訪問金鑰透過簡訊、郵件或列印紙質件交給對方，並提醒對方妥善儲存。
  \item 可新增多個聯絡人作為備份，避免單點失效。
\end{enumerate}

被指定人無需Apple裝置，憑金鑰和死亡證明即可在蘋果官網申請訪問。整個流程通常在7天內完成。

\booksubsection{Google“不活躍帳戶管理器”：自動處理身後事}

Google提供“不活躍帳戶管理器（Inactive Account Manager）”，能在使用者長期未登入時自動執行預設操作——這是目前最完善的數字遺產機制之一。

\begin{enumerate}
  \item 訪問\textbf{myaccount.google.com/inactive}，登入Google帳號。
  \item 設定超時期限：3 / 6 / 12 / 18個月未登入後觸發。
  \item 新增受信任聯絡人（最多10人），並設定每位可訪問哪些資料（Gmail、相簿、雲端硬碟等）。
  \item 可選：設定自動回覆郵件，告知來信者你已無法處理。
  \item 最後選擇：到期後是\textbf{分享資料給聯絡人}，還是\textbf{直接刪除帳戶}。
  \item 設定手機號作為兜底驗證，避免誤觸發。
\end{enumerate}

\booksubsection{國內社交帳號：紀念化與登出流程}

國內主流社交平臺對逝者帳號的處理分兩種思路：\textbf{“紀念化”（保留帳號讓親友緬懷）}和\textbf{“登出”（徹底刪除）}。家屬可憑死亡證明、親屬關係證明、申請人身份證向平臺客服申請。

\begin{itemize}
  \item \textbf{微博}：支援帳號"紀念化"，由親屬提供死亡證明申請，帳號會被鎖定無法登入但內容保留，主頁會顯示"該帳號已設定為紀念帳號"。也可申請登出，需提供繼承公證。
  \item \textbf{抖音}：透過"我 - 設定 - 反饋與幫助 - 帳號與安全 - 帳號登出"或人工客服，提交死亡證明、親屬關係證明、申請人身份證，稽核透過後可登出或轉為紀念狀態。
  \item \textbf{B站}：客服渠道提交逝者證明及親屬證明，可申請帳號登出或保留為紀念頁；UP主的歷史稿件處理需結合版權宣告。
  \item \textbf{小紅書}：透過官方客服申訴，提交死亡證明、親屬關係證明，可申請登出帳號並刪除內容。
\end{itemize}

\begin{tcbinfo}{關鍵提醒}
\begin{itemize}
  \item 各平臺規則會更新，辦理前先查閱官方客服最新指引。
  \item 所有證明檔案建議先掃描備份再提交，原件妥善保管。
  \item 跨平臺一次性辦理時，可統一申請一份\textbf{繼承公證書}，省去反覆跑手續。
\end{itemize}

\end{tcbinfo}

\booksubsection{數字資產的繼承難題}

除了社交帳號，越來越多的人持有\textbf{真正有價值的數字資產}，其繼承問題更加複雜：

\begin{itemize}
  \item \textbf{加密貨幣（比特幣、以太坊等）}：基於"私鑰即所有權"原則，\textbf{沒有私鑰就等於資產消失}。繼承人無法透過法律檔案向鏈上"申訴"。建議：將助記詞/私鑰分片存放在多個物理安全位置（如銀行保險箱+家裡保險櫃），並在遺囑中註明訪問方式。
  \item \textbf{交易所帳戶}：和銀行帳戶類似，可憑繼承公證書向交易所（如幣安、OKX）申請處置，但跨境交易所流程漫長。
  \item \textbf{域名}：屬於可繼承財產，繼承人在域名註冊商（如阿里雲、GoDaddy）提交繼承公證後可辦理過戶。
  \item \textbf{智慧財產權（原創文章、影片、專利、軟體著作權）}：著作權中的財產權利可繼承（人身權不繼承），收益權可由繼承人享有至作者死後50年。
  \item \textbf{雲端儲存與網盤}：內容資料歸屬個人，可繼承；但部分平臺協議規定帳號本身不可轉讓，需先下載再登出。
\end{itemize}

\textbf{加密貨幣“消失”的真實代價}\par

據Chainalysis報告，全球已有\textbf{超過300萬枚比特幣}因私鑰丟失而永遠無法找回，相當於現有總量的六分之一。如果你持有加密資產，請務必讓至少一位你絕對信任的人知道\textbf{恢復路徑}——否則這些財富將永遠沉睡在區塊鏈上。

\booksubsection{一份“數字遺囑”模板}

傳統的財產遺囑管不了數字資產。建議你每年更新一份“數字遺囑”附件，與正式遺囑一併存放。以下模板供參考：

\textbf{《個人數字資產處置附件》（模板）}\par

\textbf{第一部分：帳號清單}（僅記錄帳號名與所屬平臺，\textbf{不寫密碼}）

\begin{itemize}
  \item 社交類：微信（帳號xxx）、微博、抖音、B站、小紅書等
  \item 支付類：支付寶、微信支付、各銀行卡App
  \item 雲端類：iCloud、華為雲、小米雲、Google帳號、微軟帳號
  \item 郵箱類：QQ郵箱、網易郵箱、Gmail等
  \item 其他：網盤、訂閱服務、遊戲帳號等
\end{itemize}

\textbf{第二部分：密碼託管方案}

\begin{itemize}
  \item 主密碼存放方式：使用密碼管理器（如Bitwarden、1Password），主密碼寫在紙質件中，封存於家中保險櫃。
  \item 恢復金鑰/助記詞：分片存放，至少兩處物理位置（如銀行保管箱+家中保險櫃）。
  \item 授權查詢人：姓名、身份證號、與本人關係、聯絡方式。
\end{itemize}

\textbf{第三部分：處置意願}

\begin{itemize}
  \item 各類帳號希望"紀念化"還是"登出"。
  \item iCloud / Google 資料是否指定"遺產聯絡人"或"不活躍帳戶管理器"。
  \item 加密資產由誰繼承，私鑰獲取路徑說明。
  \item 版權類作品收益由誰享有。
\end{itemize}

\textbf{第四部分：聯絡人與平臺}

\begin{itemize}
  \item 信任執行人姓名、聯絡方式、備用聯絡人。
  \item 常駐平臺客服入口與所需證明材料清單。
  \item 本附件更新日期：\_\_\_\_年\_\_月\_\_日，下次複核日期：\_\_\_\_年\_\_月\_\_日。
\end{itemize}

\begin{tcbwarn}{關於"密碼"的法律提示}
將密碼直接寫進遺囑存在洩露風險，且遺囑公示後所有繼承人都能看到。\textbf{更穩妥的做法是：遺囑裡只寫“帳號清單 + 處置意願 + 密碼託管位置”，真正的密碼透過密碼管理器主密碼單獨傳給信任執行人}。這樣既符合法律要求，也降低了密碼洩露風險。

\end{tcbwarn}

\booksection{三、環保的最後一程}

一部手機裡含金、銀、銅、稀土等數十種金屬，隨意丟棄既汙染土壤，也浪費資源。正確的回收渠道有三類，各有側重：

\begin{center}
\begin{tabularx}{\linewidth}{|X|X|X|X|}
\hline
\textbf{渠道} & \textbf{回收價} & \textbf{資料安全} & \textbf{適用場景} \tabularnewline
\hline
官方以舊換新（華為/蘋果/小米等） & 中等，常疊加補貼券 & 最高，官方有規範流程 & 正打算換同品牌新機時最劃算 \tabularnewline
第三方回收平臺（愛回收、轉轉、回收寶等） & 略高，可比價 & 較高，有資料清除報告 & 想直接變現、不繫結品牌 \tabularnewline
公益捐贈（如鄉村學校、愛心機構） & 無經濟回報 & 需自行徹底擦除 & 裝置仍可用、想做善事 \tabularnewline
發揮餘熱（自留改造） & 省下新裝置錢 & 原機主本人繼續用，無風險 & 裝置效能尚可、家裡有改造需求 \tabularnewline
\hline
\end{tabularx}
\end{center}

\begin{itemize}
  \item \textbf{官方回收/以舊換新}：最推薦的方式，流程規範，資料安全有保障。蘋果、華為、小米等品牌官網都有"以舊換新"入口，評估後可抵扣新機款。
  \item \textbf{正規回收平臺}：選擇信譽良好的大型回收平臺，下單後順豐上門取件，檢測後打款。務必選擇\textbf{提供"資料清除報告"的}平臺，並保留憑證。
  \item \textbf{公益捐贈}：可聯絡當地鄉村學校、社群公益組織，舊手機可改裝成教學平板或老人機使用。\textbf{捐贈前務必執行完整的擦除流程，並卸除所有SIM/SD卡}。
  \item \textbf{發揮餘熱}：改成車載導航儀、智慧家居中控、監控攝像頭、寶寶監視器、電子相框。一臺舊手機可以再戰好幾年。
\end{itemize}

\begin{tcbinfo}{電子垃圾的真相}
一部手機的電路板裡含有\textbf{金、銀、鈀、銅}等貴金屬，每噸廢舊電路板含金量是金礦石的數十倍。但若隨意丟棄，其中的鉛、汞、鎘會滲入土壤與地下水，汙染持續數十年。\textbf{正確回收，既保護環境，也能讓稀缺金屬迴圈再生。}

\end{tcbinfo}

\begin{tcbpractice}{實操：為你的舊手機執行一場"標準葬禮"}
\begin{enumerate}
  \item 最後一次完整備份：照片上傳雲端、重要檔案複製到電腦、聊天記錄匯出存檔。
  \item 退出所有帳號：iCloud / 華為 / 小米 / Google 帳號，關閉"查詢我的手機"。
  \item 取出SIM卡與SD卡，斷開所有OTG外設。
  \item 執行加密恢復出廠設定，必要時手動勾選"格式化儲存"。
  \item （僅老裝置）用垃圾檔案填滿儲存，再恢復出廠設定一次。
  \item 選擇它的下一站：官方以舊換新、正規回收平臺、公益捐贈，或轉贈家人。
  \item 同步更新你的"數字遺產清單"，把已淘汰的裝置從清單中劃掉。
\end{enumerate}

\end{tcbpractice}

\begin{tcbpractice}{本月任務清單}
\begin{enumerate}
  \item 為你目前的主力手機指定Apple"遺產聯絡人"或設定Google"不活躍帳戶管理器"。
  \item 用本節模板起草一份《個人數字資產處置附件》，與家人溝通存放位置。
  \item 檢查家中是否有2臺以上閒置舊手機，預約一次官方或正規平臺回收。
\end{enumerate}

\end{tcbpractice}
